\documentclass[rnaas]{aastex7}

\shorttitle{SpectralUnmix}
\shortauthors{de Souza}

\begin{document}

\title{SpectralUnmix: An R Package for Smooth Spectral Unmixing of IFU and Hyperspectral Data}

\author[0000-0001-7207-4584]{Rafael S. de Souza}
\email[show]{r.s.desouza@herts.ac.uk}
\affiliation{University of Hertfordshire, Hatfield, AL10 9AB, UK}

\begin{abstract}
We present \texttt{SpectralUnmix}, an \texttt{R} package for smooth
non-negative spectral unmixing of hyperspectral data and astronomical
integral-field spectroscopy cubes. The package is designed for applications in
which a spectral cube can be approximated by a small number of non-negative
component spectra and corresponding spatial abundance maps. The current release
provides a \texttt{torch}-based optimizer, helpers for reshaping spectral cubes
into spaxel-by-wavelength matrices, plotting utilities for component maps and
component spectra, and a synthetic IFU demonstration for reproducible examples.
The package is intended as a compact and extensible software base for
astronomy-oriented source separation problems in IFU analysis, and a public
record of this initial release may be useful while more extensive validation
and benchmark studies are developed.
\end{abstract}

\section{Introduction}

Low-rank decompositions are widely used in astronomical data analysis for
compression, denoising, and exploratory analysis. In the case of
integral-field-unit (IFU) spectroscopy and other hyperspectral observations,
one often seeks a representation in which each spaxel spectrum is written as a
combination of a small number of latent components. Principal-component methods
are often useful for data compression, but their orthogonality constraints and
signed components can make direct physical interpretation difficult for some
astronomical applications. Non-negative matrix factorization (NMF) provides an
alternative factorization in which both spectral components and spatial weights
remain non-negative \citep{Lee1999,Cichocki2009}.

\texttt{SpectralUnmix} is an \texttt{R} package developed to support this style
of decomposition for IFU and hyperspectral data products. The package aims to
provide a simple software interface for reshaping cubes, fitting a smooth NMF
model, reconstructing the fitted data, and visualizing component spectra and
spatial maps. This note is intended as a citable record of the initial package
release and of the current design choices. It is not yet a full software paper:
benchmarking against alternative methods and extended real-data validation are
left to future work.

\section{Model and Package Scope}

The package assumes a linear mixing model,
\begin{equation}
\mathbf{X} \approx \mathbf{A}\mathbf{S},
\end{equation}
where $\mathbf{X}$ is a spaxel-by-wavelength matrix, $\mathbf{A}$ is a
non-negative abundance matrix, and $\mathbf{S}$ contains non-negative component
spectra. In the present implementation, optimization is performed with
\texttt{torch} \citep{Falbel2023} and includes an optional first-difference
smoothness penalty on the recovered spectra. This regularization is useful when
the desired components are expected to vary smoothly with wavelength apart from
localized spectral features.

The current package release provides:
\begin{enumerate}
\item \texttt{spectral\_unmix()} for model fitting;
\item \texttt{cube\_to\_matrix()} and \texttt{matrix\_to\_cube()} for reshaping
data products;
\item \texttt{component\_map()}, \texttt{component\_spectrum()}, and
\texttt{component\_reconstruction()} for inspecting fitted components; and
\item plotting helpers for component spectra, spatial maps, optimization loss,
and data-versus-reconstruction diagnostics.
\end{enumerate}

To support documentation and testing, the package also includes a synthetic IFU
generator that produces a small cube with continuum-like and emission-like
components. This synthetic example is not intended as an astrophysical model of
any specific source. Its purpose is to provide a fully reproducible example for
the package website, tutorials, and method development.

\section{Current Use Case}

The intended use case is straightforward. A user begins with a cube of shape
$N_x \times N_y \times N_\lambda$, reshapes it to a two-dimensional matrix,
fits a model with a user-specified number of components, and then transforms
the abundance vectors back into component maps. The fitted object also stores
the full reconstruction, which supports residual inspection and direct
comparison between observed and reconstructed spectra for selected spaxels.

For astronomy-oriented workflows, the immediate target applications are IFU
cubes from public surveys and instruments for which the dominant structures can
plausibly be described by a small number of interpretable latent spectra and
spatial distributions. Candidate future examples include MaNGA-like galaxy
cubes and MUSE observations. The current release does not yet include a
package-native FITS reader or benchmark scripts, and these are the next
priority additions before a longer journal article is attempted.

\section{Why This Note Rather Than a Full Software Paper}

The present repository contains a coherent first software release, but it does
not yet support the stronger claims usually expected from a full
\emph{Astronomy and Computing} paper. In particular, the package still needs
benchmark comparisons against standard NMF and PCA-style baselines, a public
real-data analysis vignette, and a more complete discussion of model selection
and uncertainty. A short research note is therefore the appropriate venue for
announcing the package and establishing a citable software record while the
more extensive validation work is in progress.

\section{Availability}

The source code, documentation, and manuscript materials are available in the
public \texttt{SpectralUnmix} repository:
\url{https://github.com/RafaelSdeSouza/SpectralUnmix}.

\acknowledgments
This manuscript was prepared as an initial software note for the public release
of \texttt{SpectralUnmix}. The package currently uses \texttt{torch} through
the \texttt{R} interface for tensor operations and optimization.

\bibliography{spectralunmix}

\end{document}
